\chapter{Introduzione}

\section{Obiettivi}
In questo lavoro si procederà ad analizzare la struttura dati Hash Table, 
studiandone le differenze prestazionali al variare del tipo di dato, delle 
tecniche di gestione dei conflitti e delle strategie di resizing.

Viene proposta una possibile implementazione in \texttt{C++}, controllata 
attraverso un test di correttezza presentato più avanti, che verrà utilizzata 
come base per analizzare quattro tecniche di gestione delle collisioni: 
\texttt{OPEN\_ADDRESSING\_LINEAR\_PROBING}, 
\texttt{OPEN\_ADDRESSING\_DOUBLE\_HASHING},\break
\texttt{CHAINING} e \texttt{CUCKOO}.

Verranno inoltre considerate tre strategie di resizing: 
\texttt{DOUBLE}, \texttt{CONSTANT} e \texttt{DYNAMIC}, 
che saranno spiegate nei capitoli successivi.

\section{Struttura del documento}
Il documento del progetto è costituito dai seguenti contenuti:
\begin{itemize}
    \item Analisi e descrizione dell'algoritmo Hash Table.
    \item Presentazione di una possibile implementazione in \texttt{C++}.
    \item Test di correttezza per verificare l'implementazione in \texttt{C++}.
    \item Creazione dell'apparato di testing e realizzazione di grafici per la successiva analisi.
    \item Elaborazione e analisi dei risultati ottenuti attraverso il testing.
    \item Conclusioni.
    \item Uso LLM.
\end{itemize}